\DocumentMetadata{
  pdfversion=2.0,
  pdfstandard=ua-2,
  lang=en-US,
  tagging=on,
  tagging-setup={math/setup=mathml-SE,tabsorder=structure}
}
\documentclass[11pt,letterpaper]{article}
\usepackage[margin=0.68in,includefoot]{geometry}
\usepackage{newcomputermodern}
\usepackage{amsmath,amscd,mathtools}
\usepackage{tikz,adjustbox}
\providecommand{\square}{\mdlgwhtsquare}
\usepackage[hidelinks]{hyperref}
\hypersetup{
  pdftitle={Numerical Linear Algebra First-Year Exam, January 2019},
  pdfauthor={University of Florida Department of Mathematics},
  pdfsubject={Graduate mathematics examination},
  pdfdisplaydoctitle=true,
  bookmarksopen=true
}
\setcounter{secnumdepth}{0}
\setlength{\parindent}{0pt}
\setlength{\parskip}{0.28em}
\setlength{\emergencystretch}{3em}
\setlength{\leftmargini}{2.15em}
\setlength{\leftmarginii}{2.65em}
\setlength{\leftmarginiii}{3em}
\renewcommand{\labelenumi}{\textbf{\theenumi.}}
\pagestyle{empty}
\raggedbottom
\allowdisplaybreaks
\newenvironment{examproblems}[1]{%
  \begin{enumerate}%
  \setcounter{enumi}{#1}%
  \setlength{\topsep}{0.35em}%
  \setlength{\partopsep}{0pt}%
  \setlength{\itemsep}{0.48em}%
  \setlength{\parsep}{0.18em}%
}{\end{enumerate}}
\newenvironment{examsubparts}{%
  \begin{enumerate}%
  \setlength{\topsep}{0.18em}%
  \setlength{\partopsep}{0pt}%
  \setlength{\itemsep}{0.16em}%
  \setlength{\parsep}{0.1em}%
}{\end{enumerate}}
\newenvironment{examinstructions}{%
  \par\begingroup\itshape
}{%
  \par\endgroup\vspace{0.18em}
}
\makeatletter
\renewcommand\section{\@startsection{section}{1}{0pt}%
  {-1.25ex plus -.25ex minus -.1ex}%
  {0.55ex plus .1ex}%
  {\normalfont\large\bfseries}}
\renewcommand{\@maketitle}{%
  \newpage\null\vskip -1em%
  {\footnotesize\bfseries
    \href{https://www.ufl.edu/}{University of Florida}, \href{https://math.ufl.edu/}{Department of Mathematics}\par}%
  \vskip -0.5em%
  \tagstructbegin{tag=Title}%
    {\LARGE\bfseries\href{https://gma.math.ufl.edu/exams/numerical-linear-algebra-first-year/}{Numerical Linear Algebra First-Year Exam}, January 2019\par}%
  \tagstructend%
  \vskip 0.25em%
  \tagmcbegin{artifact}%
    \hrule height 1pt%
  \tagmcend%
  \vskip 1em%
}
\makeatother
\title{Numerical Linear Algebra First-Year Exam, January 2019}
\author{}
\date{}
\begin{document}
\maketitle
\thispagestyle{empty}
\begin{examinstructions}
Do 4 (four) problems.
\end{examinstructions}
\begin{examproblems}{0}
\item
Let \(A=U\Sigma V^*\) be the singular value decomposition of \(A\in\mathbb{C}^{m\times n}\) with \(\operatorname{rank}(A)=p\) and \(p\leq n\leq m\).
\begin{examsubparts}
\item[\textnormal{(a)}]
Show \(\operatorname{Col}(A)=\operatorname{Span}\{u_1,u_2,\ldots,u_p\}\), where \(u_1,\ldots,u_p\) are the first~\(p\) columns of~\(U\).
\item[\textnormal{(b)}]
Show \(\operatorname{Null}(A^*)=\operatorname{Span}\{u_{p+1},u_{p+2},\ldots,u_m\}\).
\item[\textnormal{(c)}]
Show that \(A^*A\) is invertible if and only if~\(A\) is full rank.
\end{examsubparts}
\item
Let matrix \(A\in\mathbb{C}^{m\times n}\), with \(n<m\). Let vector \(b\in\mathbb{C}^m\), and let~\(r\) denote the residual vector \(r=b-Ax\).
\begin{examsubparts}
\item[\textnormal{(a)}]
Show that~\(x\) solves the least-squares problem \(\min \lVert b-Ax\rVert_2\) if and only if \(r\in\operatorname{Null}(A^*)\).
\item[\textnormal{(b)}]
Suppose~\(A\) is full rank, and describe how to find the least-squares solution using the QR decomposition of~\(A\).
\end{examsubparts}
\item
This problem has three parts.
\begin{examsubparts}
\item[\textnormal{(a)}]
Show that if \(A\in\mathbb{C}^{m\times m}\), \(A\) has a Schur decomposition.
\item[\textnormal{(b)}]
Show that if \(T\in\mathbb{C}^{m\times m}\) is normal and triangular, then~\(T\) is diagonal.
\item[\textnormal{(c)}]
Show that if \(A\in\mathbb{C}^{m\times m}\) is normal and~\(\lambda\) is an eigenvalue of~\(A\), then the geometric multiplicity of~\(\lambda\) is equal to the algebraic multiplicity of~\(\lambda\).
\end{examsubparts}
\item
Let \(\lVert\cdot\rVert\) be a subordinate (induced) matrix norm. If~\(A\) is \(n\times n\) invertible and~\(E\) is \(n\times n\) with \(\lVert A^{-1}\rVert\lVert E\rVert<1\), then show
\begin{examsubparts}
\item[\textnormal{(a)}]
\(A+E\) is nonsingular.
\item[\textnormal{(b)}]
\[
\lVert(A+E)^{-1}\rVert\leq \frac{\lVert A^{-1}\rVert}{1-\lVert A^{-1}\rVert\lVert E\rVert}.
\]
\end{examsubparts}
\item
Consider the matrix~\(A\) given by 
\[
A=\begin{pmatrix}1&1&2&-3\\1&10&-1&1\\2&-1&13&-2\\-3&1&-2&25\end{pmatrix}.
\]
\begin{examsubparts}
\item[\textnormal{(a)}]
Show that~\(A\) is positive definite.
\item[\textnormal{(b)}]
What can you say about the location of each of the eigenvalues of~\(A\)? Your answer should be in the form of an interval or a union of intervals.
\item[\textnormal{(c)}]
Suppose the eigenvalues of~\(A\) are all distinct (they are) and satisfy \(\lambda_1>\lambda_2>\lambda_3>\lambda_4\). Describe an algorithm that could be assured to converge to \(\lambda_4\).
\end{examsubparts}
\end{examproblems}
\end{document}
